\documentclass[11pt]{article}
\usepackage[a4paper,margin=27mm]{geometry}
\usepackage[T1]{fontenc}
\usepackage[utf8]{inputenc}
\usepackage{lmodern}
\usepackage[english]{babel}
\usepackage{microtype}
\usepackage{amsmath,amssymb,mathtools}
\usepackage{enumitem}
\pagestyle{empty}
\setlength{\parindent}{0pt}
\setlength{\parskip}{0.58em}
\setlist[enumerate]{leftmargin=2.1em,itemsep=0.35em,topsep=0.25em}
\newcommand{\OO}{\mathcal O}
\newcommand{\CC}{\mathbf C}
\newcommand{\PP}{\mathbf P}
\newcommand{\HH}{\mathcal H}
\newcommand{\uOmega}{\underline{\Omega}}
\newcommand{\eps}{\varepsilon}
\newcommand{\red}{\mathrm{red}}
\newcommand{\coh}{\mathrm{coh}}
\newcommand{\Spec}{\operatorname{Spec}}
\begin{document}
\thispagestyle{empty}

\noindent\hfill Bures, 9 October 1973

\vspace{1.1em}

Dear Illusie,

Let \(X\) be a complex analytic space, and let
\[
\eps\colon X_\ast\longrightarrow X
\]
be a smooth simplicial analytic space which is a proper hypercovering of \(X\).

\textbf{Conjecture.} The object
\[
R\eps_{\ast\ast}\bigl(\Omega^p_{X_\ast}\bigr)\in D^+_{\coh}(X,\OO)
\]
is independent of the choice of \(X_\ast\). Notation: \(\uOmega^p_X\).

\textbf{Remark.} This greatly clarifies what I do in Hodge III. One has put \(X\) into a double complex, with \(d'_1\) linear and \(d''_2\) a differential operator of first order, and well defined up to filtered quasi-isomorphism; it is \(R\eps_{\ast\ast}\Omega^\ast_{X_\ast}\), whose graded pieces are the \(R\eps_{\ast\ast}\Omega^p_{X_\ast}\).

For \(X\) projective, the spectral sequence, degenerate at \(E_1\), which abuts to the Hodge filtration, is
\[
H^q\bigl(X,\uOmega^p_X\bigr)\Longrightarrow H^{p+q}\bigl(X^{\mathrm{an}},\CC\bigr).
\]

\textbf{Evidence.} By dévissage, it should no doubt suffice to see the following. Let \(X\) be smooth, let
\(\widetilde X\xrightarrow{\pi}X\) be the blow-up of a smooth \(Y\subset X\) of codimension \(c\), and let \(\widetilde Y=\pi^{-1}(Y)\). Consider the system
\[
(\widetilde X/X)^{\Delta_n}\xrightarrow{\eps_n}X,
\]
where
\[
(\widetilde X/X)^k=\widetilde X\quad\text{glued along }\widetilde Y
\text{ to }(\widetilde Y/Y)^n
\]
\(\bigl(\)the diagonal inclusion of \(\widetilde Y\)\(\bigr)\). Let
\(\uOmega^p_{(\widetilde X/X)^k}\) be defined as
\[
\ker\Bigl(\Omega^p_{\widetilde X}\oplus
\Omega^p_{(\widetilde Y/Y)^n}\longrightarrow \Omega^p_{\widetilde Y}\Bigr).
\]
Then
\[
\Omega^p_X\xrightarrow{\sim}
R\eps_{\ast\ast}\uOmega^p_{(\widetilde X/X)^k}.
\]

``Proof''. a) Let
\[
q\colon (\widetilde Y/Y)^{\Delta_k}\longrightarrow Y.
\]
One has
\[
Rq_{\ast\ast}\Omega^p=\Omega^p
\]
\(\bigl(\)clear: there are local sections\(\bigr)\).

b) This makes it possible to simplify the diagram. The geometric fact should be that \(X\) is the geometric realization of
\[
\widetilde X\rightrightarrows \widetilde X\times_X\widetilde X,
\]
and one must verify the triangle
\[
0\longrightarrow \Omega^p_X\longrightarrow
\Omega^p_Y\oplus R\pi_\ast\Omega^p_{\widetilde X}
\longrightarrow R\pi'_\ast\Omega^p_{\widetilde Y}\longrightarrow 0.
\]
That is:
\[
\Omega^p_X\simeq \pi_\ast\Omega^p_{\widetilde X},\qquad
\Omega^p_Y\simeq \pi'_\ast\Omega^p_{\widetilde Y},
\]
and, for \(i>0\),
\[
R^i\pi_\ast\Omega^p_{\widetilde X}\simeq
R^i\pi'_\ast\Omega^p_{\widetilde Y}.
\]
Indeed, one has
\[
0\longrightarrow \pi^\ast\Omega^1_X\longrightarrow
\Omega^1_{\widetilde X}\longrightarrow \Omega^1_{\widetilde Y/Y}\longrightarrow 0,
\]
whence a filtration of \(\Omega^p_{\widetilde X}\) with successive quotients
\[
\Omega^p_{\widetilde Y/Y},\quad \ldots,\quad
\Omega^p_{\widetilde Y/Y}(i)\ (0<i<p),\quad\text{and}\quad
\pi^\ast\Omega^p_X,
\]
the former having \(R\pi_\ast=0\), and the latter having \(R\pi_\ast=\Omega^p_X\). The assertion follows.

\textbf{Example.} For \(X\) projective, if \(\uOmega^0_X=\OO\), then
\[
H^\ast(X^{\mathrm{an}},\CC)\longrightarrow H^\ast(X,\OO)
\]
is surjective.

\textbf{Example.} \(\HH^0(\uOmega^0_X)\) is the subsheaf, on \(X\), of the sheaf of rings of the normalization of \(X\), formed by the functions that are constant on each reduced fiber. Denoting the normalization by \(n\), one has
\[
0\longrightarrow \HH^0(\uOmega^0_X)\longrightarrow
\OO_{X^n}\rightrightarrows
\OO_{(X^n\times_X X^n)_{\red}}.
\]

\textbf{Example.} One should have, for a finite group \(G\) acting on \(X\), and \(\pi\colon X\to X/G\),
\[
\uOmega_{X/G}=\bigl(\pi_\ast\uOmega_X\bigr)^G.
\]
Where \(Z=\) smooth/finite group, one should therefore have \(\uOmega^0=\OO\). For a proper V-variety \(X\), locally \(X=X_1/G\), and
\[
\uOmega^p_X=\bigl(\pi_\ast\Omega^p_{X_1}\bigr)^G,
\]
the complex \(\uOmega^\bullet_X\) then gives rise to a standard Hodge theory.

\textbf{Example.} If \(A=B\cup_C D\), one should have a triangle
\[
0\longrightarrow \uOmega_A\longrightarrow
\uOmega_B\oplus\uOmega_D\longrightarrow \uOmega_C\longrightarrow 0.
\]
If \(\uOmega^0=\OO\) for \(B\), \(D\), and \(C\), it should therefore be the same for \(A\).

\textbf{Example.} If \(\widetilde X\xrightarrow{\pi}X\) is birational, and an isomorphism outside \(Y\subset X\), with inverse image \(\widetilde Y\), \(\pi'\colon\widetilde Y\to Y\), one should have a triangle
\[
0\longrightarrow \uOmega_X\longrightarrow
\uOmega_Y\oplus R\pi_\ast\uOmega_{\widetilde X}
\longrightarrow R\pi'_\ast\uOmega_{\widetilde Y}\longrightarrow 0.
\]

\textbf{Example.} Let \(X\) be the cone over a projective variety \(V\subset\PP^e\). One has:
\[
\Gamma\HH^0(\uOmega^0_X)=\CC\oplus\bigoplus_{n>0}H^0\bigl(V,\OO(n)\bigr),
\]
and \(\Gamma\HH^1(\uOmega^0_X)\) is the following graded module over this graded algebra:
\[
\bigoplus_{n>0}H^1\bigl(V,\OO(n)\bigr).
\]
One has \(\uOmega^0_X=\OO\) if and only if, for \(n>0\),
\[
H^\ast\bigl(\PP^e,\OO(n)\bigr)\longrightarrow H^\ast\bigl(V,\OO(n)\bigr)
\]
is surjective. Example: \(V=V_m(\underline a)\), \(\sum a_i<n+2\) \((\mathrm{FAC}\ 78,\ \mathrm{Prop.}\ 5)\).

\textbf{Functoriality.}
\begin{enumerate}[label=\alph*)]
\item \(\uOmega_X\) should be functorial in \(X\).
\item There is a ``cup product'', at least:
\[
\HH^i\uOmega^p\otimes\HH^j\uOmega^q\longrightarrow
\HH^{i+j}\uOmega^{p+q}.
\]
\item If \(X_\ast\) is a proper hypercovering of \(X\), one should have
\[
R\eps_{\ast\ast}\Omega^p_{X_\ast}=\uOmega^p_X.
\]
\end{enumerate}

\textbf{Generalization.}
\begin{enumerate}[label=\alph*)]
\item Ideally, one should have such a theory for sheaves other than \(\CC\) \(\bigl(\)complexes of sheaves\(\bigr)\).
\item At least: \(U\hookrightarrow X\), \(Rj_\ast\CC\) \((\)cf. Hodge III\()\).
\end{enumerate}

I hope these scribbles are not too garbled.

Best regards,

\hspace{2em}P. Deligne

\end{document}
